\scanpage{264}\mspage{180}
% archive sheet number: 263
En effet, la condition a) des \unsure{fcg|feg} (cf 4.7.7) résulte de
la $=$ condition pour $(\mathcal F^{\unsure{\pi}})$, et de (4.7.10.1)
$M_\pi\subset\mathcal F^{\pi^{\pi_0}}$. La condition b) est une \deleted{\unsure{condition}} \inserted{\unsure{propriété}}
générale de la $\pi$-accessibilité, compte tenu
que $(\mathcal F^{\pi_0})_0$ est une petite sous-catégorie génératrice
formée d'objets $\pi_0$-accessibles de $M$. La condition \unsure{d|b}) est satisfaite tautologiquement, et
il reste à vérifier la condition c). Mais
on a vu que pour $\pi=\pi^{\pi_0}$, on a $M_\pi=\mathcal F^\pi$
(4.7.9.1), donc la condition c) pour $(M_\pi)$
résulte de la $=$ condition pour $(\mathcal F^\pi)$.

Ainsi, on trouve

\uline{Proposition} 4.7.15. Si $(\mathcal F^\pi)_{\pi\geq\unsure{\pi_0|\pi}}$ est une
fcga sur la $\unsure{U|\mathcal U}$-catégorie $M$, il en est
de $=$ de $(M_\pi)_{\pi\geq\pi_0}$. De plus \deleted{tout}
cardinal utile (resp. régulier) pour la
première filtration, l'est pour la \deleted{\unsure{f}} seconde
(appelée la \deleted{filtration cardinale} \inserted{fcga \uline{canonique}} \uline{associée}
à la filtration cardinale donnée).

La dernière assertion résulte du
fait que si $\pi$ est utile, pour $(\mathcal F^\pi)$, alors
on a $\mathcal F^\pi=M_\pi$..

4.7.16. Il me semble qu'on \unsure{peut} dire que la
fcg canonique associée à une fcga donnée
